\section*{第 \ref{ch:b1:stress-strain-rheology} 章——应力、应变与流变学}

\begin{solution}{exo:b1:stress-strain-rheology:1}
(a) $\sigma_v = 2700\times9.81\times10\,000 = \qty{2.65e8}{Pa} = \qty{264.9}{MPa}$。
(b) $P = 1000\times9.81\times10\,000 = \qty{98.1}{MPa}$。(c) $\sigma_v' =
264.9 - 98.1 = \qty{166.8}{MPa}$，且 $\lambda = 98.1/264.9 = 0.370$，即
$1000/2700$.
\end{solution}

\begin{solution}{exo:b1:stress-strain-rheology:2}
(a) $e = (13-10)/10 = 0.30$，$S = 1.30$。(b) $e = (7-10)/10 = -0.30$，即
缩短 \qty{30}{\percent}，$S = 0.70$。(c) 该角已闭合
$\psi = \ang{30}$，且 $\gamma = \tan\ang{30} = 0.577$。
\end{solution}

\begin{solution}{exo:b1:stress-strain-rheology:3}
(a) 圆心为 \qty{70}{MPa}，半径为 \qty{30}{MPa}。(b) 最大切应力
即半径，\qty{30}{MPa}，作用于与 $\sigma_1$ 和
$\sigma_3$ 成 \ang{45} 的面上，这些面上的\omterm{def:b1:stress-strain-rheology:stress}{法向应力}为圆心处的值 \qty{70}{MPa}。(c)
$2\theta = \ang{60}$：$\sigma_n = 70 + 30\cos\ang{60} = \qty{85}{MPa}$，且
$\tau = 30\sin\ang{60} = \qty{26.0}{MPa}$.
\end{solution}

\begin{solution}{exo:b1:stress-strain-rheology:4}
(a) $\ang{45} - \ang{30}/2 = \ang{30}$，即两平面与 $\sigma_1$ 成此角，分居其两侧，
故两平面之间的夹角为 \ang{60}。(b) 令 $c' = 20 + r$，
由 \cref{thm:b1:stress-strain-rheology:failure-angle} 的条件得
$r = C\cos\phi + (20 + r)\sin\phi$，于是
\[
  r = \frac{C\cos\phi + 20\sin\phi}{1 - \sin\phi} = \frac{8.66 + 10}{0.5}
    = \qty{37.3}{MPa},
\]
且 $\sigma_1' = 20 + 2r = \qty{94.6}{MPa}$。
\end{solution}

\begin{solution}{exo:b1:stress-strain-rheology:5}
(a) $c = 40$，$R = \sqrt{20^2 + 15^2} = \qty{25}{MPa}$：$\sigma_1 =
\qty{65}{MPa}$ 及 $\sigma_3 = \qty{15}{MPa}$；$\tan2\theta_p = 30/40 = 0.75$，
$2\theta_p = \ang{36.87}$，$\theta_p = \ang{18.43}$。(b) 此角度下
$\cos^2\theta = 0.9$，$\sin^2\theta = 0.1$ 且 $\sin\theta\cos\theta = 0.3$：
$\sigma_n = 60\times0.9 + 2\times15\times0.3 + 20\times0.1 = 54 + 9 + 2 =
\qty{65}{MPa}$，且 $\tau = (20-60)\times0.3 + 15\times(0.9-0.1) = -12 + 12 =
0$。(c) $R = \qty{25}{MPa}$，作用于
与主方向成 \ang{45} 的面上。
\end{solution}

\begin{solution}{exo:b1:stress-strain-rheology:6}
(a) $S_{1,2} = \sqrt{1 + 1} \pm 1$：$S_1 = 2.414$，且 $S_2 = 0.414$；
$\tan2\theta' = 2/2 = 1$，$\theta' = \ang{22.5}$。(b) $S_1S_2 =
(\sqrt2 + 1)(\sqrt2 - 1) = 1$，即 $\det F$：面积不变。(c) 沿
$(1, 1)$ 的线段变为 $F(1, 1) = (1 + 2, 1) = (3, 1)$，与
剪切方向的夹角为 $\arctan(1/3) = \ang{18.4}$，\omterm{def:b1:stress-strain-rheology:strain}{伸长比}为
$\sqrt{10}/\sqrt2 = 2.24$。该线段已越过椭圆的
长轴（位于 \ang{22.5}），朝剪切面方向转动；其\omterm{def:b1:stress-strain-rheology:strain}{伸长比}略小于
$S_1$，这是必然的，因为任何线段的\omterm{def:b1:stress-strain-rheology:strain}{伸长比}都不会超过 $S_1$。
\end{solution}

\begin{solution}{exo:b1:stress-strain-rheology:7}
(a) $S_1/S_2 = 4$ 及 $S_1S_2 = 1$ 给出 $S_1 = 2$，$S_2 = 0.5$；由
\cref{prop:b1:stress-strain-rheology:simple-shear} 可知 $\gamma = S_1 - S_2 =
1.5$。(b) $\tan2\theta' = 2/1.5 = 1.333$，$2\theta' = \ang{53.1}$：长轴应
与边界成 \ang{26.6}。(c) $1.5\times200 = \qty{300}{m}$。
\end{solution}

\begin{solution}{exo:b1:stress-strain-rheology:8}
$\sigma_v = 2700\times9.81\times3000 = \qty{79.46}{MPa}$。(a) $P = \qty{29.43}{MPa}$，
$\sigma_n' = 70 - 29.43 = \qty{40.57}{MPa}$，\omterm{prop:b1:stress-strain-rheology:byerlee}{拜尔利定律}允许的值为
$0.85\times40.57 = \qty{34.5}{MPa}$，大于所施加的 \qty{30}{MPa}：断层
保持稳定。(b) $P = 0.6\times79.46 = \qty{47.68}{MPa}$，$\sigma_n' =
\qty{22.32}{MPa}$，摩擦力仅给出 \qty{19.0}{MPa}：断层
发生滑动。(c) $30 = 0.85\,(70 - 79.46\,\lambda)$ 给出 $\lambda = (70 -
35.29)/79.46 = 0.437$.
\end{solution}

\begin{solution}{exo:b1:stress-strain-rheology:9}
(a) 应力正比于 $\dot\varepsilon^{1/4}$：$10^{1/4} = 1.78$。(b) 两者
强度之比为
\[
  \exp\Bigl[\frac{E}{nR}\Bigl(\frac{1}{673.15} - \frac{1}{773.15}\Bigr)\Bigr]
  = \exp\bigl(6705\times\num{1.921e-4}\bigr) = \exp(1.288) = 3.63 .
\]
(c) 应变速率变化十倍，蠕变强度改变 1.78 倍，
而温度只需变化 $100\times\ln1.78/\ln3.63 \approx
\qty{45}{\degreeCelsius}$ 即可产生同样效果；在地温梯度约为 \qty{20}{\degreeCelsius}
每千米的\omterm{def:b1:heat-budget:geotherm}{地温线}上，这会使转换带移动约 \qty{2}{km}。\omterm{def:b1:heat-budget:geotherm}{地温线}的变化
会使等温线随之移动，转换带的移动幅度远大于此。
\end{solution}

\begin{solution}{exo:b1:stress-strain-rheology:10}
(a) 当 $C = 0$ 时，断层在 $r = \mu c'/\sqrt{1+\mu^2}$
时发生滑动（\cref{thm:b1:stress-strain-rheology:failure-angle}）。此处 $c' = S =
\rho g z(1-\lambda)$，故 $\sigma_1 - \sigma_3 = 2r = 2\mu S/\sqrt{1+\mu^2}$。
(b) $S$ 每千米增加 $2700\times9.81\times1000\times(1 - 1000/2700) =
\qty{16.68}{MPa}$，且 $2\times0.85/1.3124 = 1.295$：斜率为
\qty{21.6}{MPa} 每千米，介于 \num{13.1}（正断层）与
\qty{61.3}{MPa} 每千米（逆冲断层）之间。(c) $2\mu/(k+\mu) = 0.786$，
$2\mu/k = 1.295$ 及 $2\mu/(k-\mu) = 3.68$，取 $k = 1.3124$：滑动时的
\omterm{def:b1:stress-strain-rheology:principal}{差应力}分别为 0.79、1.30 和 3.68 倍的竖直
\omterm{def:b1:stress-strain-rheology:lithostatic}{有效应力}。
\end{solution}

\begin{solution}{exo:b1:stress-strain-rheology:11}
(a) $t = 60\times\num{e6}\times\num{3.156e7} = \qty{1.89e15}{s}$，且
$2\sqrt{\kappa t} = 2\sqrt{\num{1.89e9}} = \qty{87.0}{km}$。(b)
$\operatorname{erf}x = 600/1350 = 0.444$ 对应 $x = 0.417$：$z = 0.417\times87.0
= \qty{36.3}{km}$；$\operatorname{erf}x = 800/1350 = 0.593$ 对应 $x = 0.586$：
$z = \qty{51.0}{km}$。(c) 地震应当达到 \qtyrange{36}{51}{km}；
包络线将脆性层底部置于 \qty{38.2}{km}，位于
该范围之内且接近其顶部，对应温度为 \qty{641}{\degreeCelsius}，
与观测范围相符。
\end{solution}

\begin{solution}{exo:b1:stress-strain-rheology:12}
(a) $r = 45$ 且 $c' = \qty{65}{MPa}$；圆心到直线
$\tau = 0.85\sigma_n'$ 的距离为 $0.85\times65/1.3124 = \qty{42.1}{MPa}$，小于
半径，圆与直线相交，位于其上方圆弧上的面
发生滑动。(b) 以 $2\theta$ 表示自 $\sigma_1$ 端起算的角，滑动需满足
$45\sin2\theta \ge 0.85\,(65 + 45\cos2\theta)$，即
$45\sin2\theta - 38.25\cos2\theta \ge 55.25$，或
$59.06\sin(2\theta - \ang{40.36}) \ge 55.25$，其中 $\ang{40.36} = \phi$。
故 $\sin(2\theta - \phi) \ge 0.9355$，$2\theta - \phi$ 介于 \ang{69.31}
与 \ang{110.69} 之间，$2\theta$ 介于 \ang{109.67} 与 \ang{151.06} 之间：法线
与 $\sigma_1$ 的夹角介于 \ang{54.8} 至 \ang{75.5} 之间，而滑动
破裂面与 $\sigma_1$ 的夹角介于 \ang{14.5} 至 \ang{35.2} 之间。范围的中点
\ang{24.8}，即 $\ang{45} - \phi/2$，正是
\cref{thm:b1:stress-strain-rheology:failure-angle}所给的面。(c) 只有该面在
圆恰好与直线相切时发生滑动：$r = \mu(20 + r)/k$，$r = 20\mu/(k -
\mu) = 17/0.4624 = \qty{36.8}{MPa}$，且 $\sigma_1' = 20 + 2r =
\qty{93.5}{MPa}$.
\end{solution}

\begin{solution}{pb:b1:stress-strain-rheology:1}
\textbf{1.} $\sigma_v = 2700\times9.81\times5000 = \qty{132.44}{MPa}$.

\textbf{2.} $P = 1000\times9.81\times5000 = \qty{49.05}{MPa}$，且
$\lambda = 49.05/132.44 = 0.370$.

\textbf{3.} $\sigma_1 = \sigma_v = \qty{132.44}{MPa}$，竖直，且
$\sigma_3 = \qty{75.00}{MPa}$，水平东西向：\omterm{def:b1:layered-earth:layers}{地壳}正在
在自身重力作用下东西向拉伸，这种应力状态驱动正断层，
如本例中的断层。

\textbf{4.} $\sigma_1 - \sigma_3 = \qty{57.44}{MPa}$；圆心
$(132.44 + 75.00)/2 = \qty{103.72}{MPa}$；半径 \qty{28.72}{MPa}。

\textbf{5.} $\sigma_1' = 132.44 - 49.05 = \qty{83.39}{MPa}$, $\sigma_3' =
75.00 - 49.05 = \qty{25.95}{MPa}$，圆心 $103.72 - 49.05 = \qty{54.67}{MPa}$。

\textbf{6.} $\sigma_{xx} = \qty{75.00}{MPa}$, $\sigma_{yy} = \qty{132.44}{MPa}$
且 $\sigma_{xy} = 0$：坐标轴即为主方向，其中一个
竖直，因此与之垂直的面上不存在切应力。

\textbf{7.} 沿断层\omterm{def:b1:geological-maps:attitude}{倾向}向下，断层沿单位向量
$(\cos\ang{60}, -\sin\ang{60})$ 延伸，即向东并向下；其向上的法向量为
$\vec n = (\sin\ang{60}, \cos\ang{60}) = (0.866, 0.500)$，与
$x$ 轴夹角为 \ang{30}。其与 $\sigma_1$（沿 $y$ 方向）的夹角余弦为
0.500: $\theta = \ang{60}$.

\textbf{8.} 取与 $x$ 成 \ang{30} 角的面：$\sigma_n = 75.00\times0.75 + 132.44\times0.25
= 56.25 + 33.11 = \qty{89.36}{MPa}$，且 $\tau = (132.44 - 75.00)\times
\sin\ang{30}\cos\ang{30} = 57.44\times0.4330 = \qty{24.87}{MPa}$.

\textbf{9.} 取 $2\theta = \ang{120}$（自 $\sigma_1$ 端起算）：
$\sigma_n = 103.72 + 28.72\cos\ang{120} = 103.72 - 14.36 = \qty{89.36}{MPa}$
且 $\tau = 28.72\sin\ang{120} = \qty{24.87}{MPa}$。

\textbf{10.} $\sigma_n' = 89.36 - 49.05 = \qty{40.31}{MPa}$.

\textbf{11.} $24.87/40.31 = 0.617$，低于 0.85：断层保持稳定。

\textbf{12.} $0.85\times40.31 = \qty{34.26}{MPa}$；切应力
低于该值 \qty{9.39}{MPa}。

\textbf{13.} 圆心位于 \qty{54.67}{MPa}，半径
\qty{28.72}{MPa}，圆的范围从 \num{25.95} 到 \qty{83.39}{MPa}；拜尔利直线为
$\tau = 0.85\,\sigma_n'$（此处所有正应力均低于
\qty{200}{MPa}）；断层即点 $(40.31, 24.87)$，
在自 $\sigma_1$ 端起算 \ang{120} 处。

\textbf{14.} $0.85\times54.67/\sqrt{1 + 0.85^2} = 46.47/1.3124 =
\qty{35.41}{MPa}$，大于半径 \qty{28.72}{MPa}：整个圆
位于直线下方，岩石中任何面都不会滑动。

\textbf{15.} $\phi = \arctan0.85 = \ang{40.36}$：该面位于
$\ang{45} - \phi/2 = \ang{24.8}$ 处，与竖直的 $\sigma_1$ 成，\omterm{def:b1:geological-maps:attitude}{倾角}为
\ang{65.2}.

\textbf{16.} $S = \sigma_v' = \qty{83.39}{MPa}$，且
$\sigma_1 - \sigma_3 = 2\times0.85\times83.39/(1.3124 + 0.85) =
\qty{65.55}{MPa}$；由斜率得 $13.11\times5 = \qty{65.55}{MPa}$。

\textbf{17.} $\sigma_3 = 132.44 - 65.55 = \qty{66.88}{MPa}$.

\textbf{18.} 断层发生滑动的条件为
$r\sin\ang{120} = 0.85\,(c' + r\cos\ang{120})$，其中
$c' = (83.39 + \sigma_3')/2$ 与 $r = (83.39 - \sigma_3')/2$；解得
$\sigma_3' = \qty{17.18}{MPa}$，$\sigma_3 = \qty{66.23}{MPa}$。该值低于
最佳平面的情形：断层并不位于圆首次与直线接触的那一点，
因此圆必须进一步增大，方能到达
断层自身所对应的点。

\textbf{19.} 所有有效正应力，包括 $\sigma_1'$ 与 $\sigma_3'$，
均减小 $\Delta P$，而二者之差保持不变；
切应力则完全不依赖孔隙压力。

\textbf{20.} 当距离等于半径时，最优面开始滑动：
$0.85\,(54.67 - \Delta P)/1.3124 = 28.72$，故
$\Delta P = 54.67 - 28.72\times1.3124/0.85 = 54.67 - 44.34 = \qty{10.33}{MPa}$.

\textbf{21.} 井中水柱的重量与原来相同，因此
泵须在井口补加全部 \qty{10.33}{MPa}，约为 103 倍
大气压强：一根水柱的重量相当于高
$\num{10.33e6}/(1000\times9.81) = \qty{1053}{m}$ 的水柱。

\textbf{22.} 取 $\mu = 0.6$，断层只能承受
$0.6\times40.31 = \qty{24.19}{MPa}$，小于其承受的
\qty{24.87}{MPa}，则断层早已滑动。既然它仍保持静止，其摩擦系数至少为
$0.617$，否则必须具有一定的\omterm{def:b1:stress-strain-rheology:coulomb}{内聚力}。

\textbf{23.} 取向最佳的面在孔隙压力达到
$49.05 + 10.33 = \qty{59.38}{MPa}$ 时开始滑动，即 $\lambda = 59.38/132.44 = 0.448$，
远低于\omterm{def:b1:stress-strain-rheology:lithostatic}{静岩压力}对应的值 1。

\textbf{24.} $\Delta P = \sigma_n' - \tau/0.85 = 40.309 - 24.870/0.85 = 40.309 -
29.259 = \qty{11.049}{MPa}$，大于能使最佳取向的面
开始滑动的 \qty{10.33}{MPa}。保留一位小数，
使断层滑动的孔隙压力升高量为
\textbf{\qty[text-series-to-math=true]{11.0}{MPa}}.
\end{solution}
